\subsubsection{From QUBO to Cost Hamiltonian}\label{sec:from-qubo-to-cost-hamiltonian}

We have:
\begin{itemize}
    \item An optimization problem;
    \item The optimization problem is expressed as a QUBO problem;
    \item The $Q$ matrix is known and contains the coefficients of the QUBO problem.
\end{itemize}
But a quantum computer cannot optimize binary variables $x_i \in \{0, 1\}$ directly. Instead, it evolves \textbf{quantum states} under a \textbf{Hamiltonian}. So in this section we will see \textbf{\emph{how to convert a classical QUBO into a quantum Cost Hamiltonian}}.

\highspace
\begin{flushleft}
    \textcolor{Green3}{\faIcon{link} \textbf{Step 1: Binary variables $\to$ Spin variables (QUBO $\to$ Ising model)}}
\end{flushleft}
A quantum computer does \textbf{not naturally understand binary variables} $x_i \in \{0, 1\}$. It naturally manipulates \textbf{qubits}, whose computational basis states are eigenstates of the Pauli-Z operator. The eigenvalues of the Pauli-Z operator are $\pm 1$, which is not the same as the binary values $0$ and $1$, but are exactly the values of \textbf{spin variables}! So we first rewrite the optimization problem using spins.

\highspace
Suppose we have a QUBO problem:
\begin{equation*}
    \min_{x} x^T Q x \quad \text{s.t. } x_i \in \{0, 1\}
\end{equation*}
Substitute the binary variables $x_i$ with spin variables $s_i$ using the relation:
\begin{equation}\label{eq:binary-to-spin}
    s = 2x - 1 \quad \Longleftrightarrow \quad x = \frac{1 + s}{2}
\end{equation}
After substituting $x_i = \frac{1 + s_i}{2}$ into the QUBO objective function (\autopageref{eq:qubo-formulation}):
\begin{equation*}
    y = \sum_{i} a_{i} x_{i} + \sum_{i > j} q_{ij} x_{i} x_{j}
\end{equation*}
Every QUBO can be rewritten as an Ising model (\autopageref{eq:ising-energy-function}):
\begin{equation*}
    \underbrace{y}_{H_{\text{Ising}}} = c + \sum_i S_{ii} s_i + \sum_{i>j} S_{ij} s_i s_j
\end{equation*}
Where:
\begin{itemize}
    \item $c$ is a constant term that does not depend on the spin variables $s_i$;
    \item $S_{ii}$ are the local fields (linear coefficients);
    \item $S_{ij}$ are the coupling strengths (quadratic coefficients).
\end{itemize}
Notice that the Ising model has the same structure as the QUBO problem: a linear term (the first sum) and a quadratic term (the second sum). The only difference is that the variables are now spins $s_i \in \{-1, +1\}$ instead of binary variables $x_i \in \{0, 1\}$. Also, the matrix $Q$ is \textbf{replaced} by a new matrix $S$ that contains the \textbf{coefficients of the Ising model}, but \textbf{describe exactly the same optimization problem} as the original QUBO.

\highspace
\begin{flushleft}
    \textcolor{Green3}{\faIcon{link} \textbf{Step 2: From Ising Model to Cost Hamiltonian}}
\end{flushleft}
The Ising model is a classical model that describes the energy of a system of spins. However, the spins are still classical objects that can only be in one of two possible states. To ``\emph{quantize}'' the Ising model, we replace the classical spin variables $s_i$ with quantum operators. The natural choice for this is the Pauli-$Z$ operator, which has eigenvalues $\pm 1$ and acts on qubits (\autopageref{sec:quantum-ising-model}). The resulting quantum Hamiltonian is called the \definition{Cost Hamiltonian}:
\begin{equation}\label{eq:cost-hamiltonian}
    H_C = \sum_i S_{ii} \sigma_{Z}^{i} + \sum_{i > j} S_{ij} \sigma_{Z}^{i} \sigma_{Z}^{j}
\end{equation}
Nothing of difficult has happened here: we have \textbf{simply replaced the classical spin variables $s_i$ with quantum operators} $\sigma_{Z}^{i}$, which act on the $i$-th qubit. This simple \textbf{replacement is allowed} because the \textbf{Pauli-Z operator preserves the mathematical structure of the optimization problem} while turning it into a quantum operator:
\begin{equation*}
    Z\ket{0} = \ket{0}, \quad Z\ket{1} = -\ket{1}, \quad \text{with eigenvalues } +1 \text{ and } -1
\end{equation*}

\highspace
The resulting Hamiltonian $H_C$ is a matrix of size $2^n \times 2^n$, where $n$ is the number of qubits. This exponential growth in size is due to the fact that each term $\sigma_{Z}^{i}$ acts on a different qubit, and the total Hilbert space is the tensor product of all the qubits ($I \otimes I \otimes \dots \otimes \sigma_Z^i \otimes \dots \otimes I$).

\highspace
The Cost Hamiltonian is a quantum operator that encodes the same optimization problem as the original QUBO, but in a form that can be manipulated by a quantum computer. Also, it is \textbf{diagonal} and \textbf{enumerates the energy of all possible states}, which is exactly what we need for optimization.

\highspace
\textcolor{Green3}{\faIcon{question-circle} \textbf{Why does the constant disappear?}} The constant term $c$ does not appear in the Hamiltonian because \hl{adding a constant to every energy level does \textbf{not} change which state has the lowest energy}. Since optimization only depends on the ordering of the energies, the constant term can be safely ignored.

\highspace
In summary, to obtain the Cost Hamiltonian from a QUBO problem, we follow these steps:
\begin{enumerate}
    \item The QUBO formulation is rewritten in terms of spin variables using the transformation $s=2x-1$. This produces a spin objective function of the form:
    \begin{equation*}
        y = c + \sum_i S_{ii} s_i + \sum_{i>j} S_{ij} s_i s_j
    \end{equation*}
    \item Each spin variable $s_i$ is replaced by the corresponding Pauli-$Z$ operator $\sigma_{Z}^{i}$, obtaining the Cost Hamiltonian:
    \begin{equation*}
        H_C = \sum_i S_{ii} \sigma_{Z}^{i} + \sum_{i > j} S_{ij} \sigma_{Z}^{i} \sigma_{Z}^{j}
    \end{equation*}
    The constant term can be discarded because it shifts all energies equally.
\end{enumerate}
The resulting Hamiltonian is diagonal, acts on quantum states, and its ground state corresponds to the optimal solution of the original optimization problem.

\highspace
At this point, we have successfully built the Cost Hamiltonian. The remaining question is purely practical: \emph{how do we implement} $U_C(\gamma) = e^{-i \gamma H_C}$ (i.e., the unitary operator corresponding to the Cost Hamiltonian) \emph{as an actual quantum circuit?} The next section answers this.